% ---------------------------------------------------------------------------
% A5ModalRuntime.tex
% Altitude A5 — Modal Runtime Physics
% Author: Borealis S. Hedling
% Version: v1.0
% ---------------------------------------------------------------------------

\section{Altitude A5: Modal Runtime Physics}

Altitude A5 introduces modal dynamics on top of the spectral runtime manifold
defined in A4. Modal runtime physics treats spectral modes as operational
entities, enabling runtime patches to evolve according to modal coordinates,
modal operators, and modal invariants.

A5 is the final altitude before dual algebra (A6) and non-dual algebra (A7).

\subsection{Modal Manifold}

Let $M$ be the spectral runtime manifold from A4. A5 defines the modal manifold:
\[
\mathcal{M}_{\mathrm{modal}} = (M, \Psi, \mathcal{T}_{\mathrm{modal}}),
\]
where:
\begin{itemize}
  \item $\Psi$ is the modal coordinate embedding,
  \item $\mathcal{T}_{\mathrm{modal}}$ is the admissible modal operator set.
\end{itemize}

Modal coordinates $\psi_k$ from A4 become active runtime variables.

\subsection{Modal Operators}

A modal operator $T_{\mathrm{modal}}$ acts on modal coordinates:
\[
T_{\mathrm{modal}} : \Psi(p) \mapsto \Psi'(p).
\]

A5 requires:
\[
\Psi'(p) = f(\psi_1(p), \ldots, \psi_d(p)),
\]
where $f$ is:
\begin{itemize}
  \item continuous,
  \item adjacency-preserving,
  \item spectrally stable,
  \item reversible.
\end{itemize}

\subsection{Modal Dynamics}

Modal dynamics define how runtime patches evolve:
\[
\frac{d\psi_k}{dt} = F_k(\psi_1, \ldots, \psi_d),
\]
where $F_k$ must satisfy:
\begin{itemize}
  \item spectral continuity,
  \item modal stability,
  \item bounded holonomy,
  \item membrane consistency.
\end{itemize}

Modal dynamics must not fracture adjacency or violate spectral invariants.

\subsection{Modal Stability}

A5 enforces:
\[
\left| \frac{d\psi_k}{dt} \right| < \epsilon,
\]
ensuring modal evolution remains compatible with A4 spectral continuity.

Modal stability prevents:
\begin{itemize}
  \item spectral tearing,
  \item modal divergence,
  \item adjacency fracture.
\end{itemize}

\subsection{Modal Reconstruction}

A5 requires reversible reconstruction:
\[
p = \Psi^{-1}(\Psi(p)).
\]

This ensures modal evolution does not lose geometric information.

\subsection{Admissible Modal Operators}

Admissible operators include:
\begin{itemize}
  \item modal filters,
  \item modal projections,
  \item modal integrators,
  \item modal evolution operators.
\end{itemize}

These operators must preserve:
\begin{itemize}
  \item adjacency,
  \item spectral continuity,
  \item modal stability,
  \item provenance monotonicity.
\end{itemize}

\subsection{DiffusionGeometry as a Modal Example}

DiffusionGeometry provides modal examples:
\begin{itemize}
  \item diffusion coordinates $\psi_k$,
  \item spectral decay modes,
  \item modal embeddings,
  \item modal evolution under diffusion time $t$.
\end{itemize}

These examples remain strictly within altitudes A3–A5.

\subsection{Integration with Python Runtime Harness (A4.5)}

The Python runtime harness enforces:
\begin{itemize}
  \item modal coordinate updates,
  \item reversible modal transitions,
  \item adjacency-preserving modal evolution,
  \item provenance monotonicity under modal dynamics.
\end{itemize}

A5ModalRuntime.tex defines the modal invariants the harness must satisfy.

% ---------------------------------------------------------------------------
% Provenance
% A5ModalRuntime.tex
% Author: Borealis S. Hedling
% Altitude: A5
% Surface: Structural-AI/Runtime-Physics
% Version: v1.0
% Created: 18 October 2026, Dublin, Ireland
% Purpose: Define modal runtime physics, modal operators, modal dynamics,
%          modal stability, reversible modal reconstruction, and admissible
%          modal evolution. Provides theoretical foundation for A6 dual algebra
%          and integrates with the Python A4.5 runtime harness.
% ---------------------------------------------------------------------------
